\section{From a Boolean marker to a custom scan}
\label{sec:mmsearch-sql-rendezvous}

% \texttt{mmsearch} is declared as an ordinary SQL function returning
% \texttt{boolean}. PostgreSQL can therefore resolve its arguments and accept
% it in a condition without a new SQL grammar. Its two-argument form uses the
% session's alignment settings. A third argument can supply typed options.
% The extra hit rows are part of BioSeqDB's search contract, which every
% accepted plan must preserve.
% During planning the call acts as a \emph{marker}: it identifies the search
% that an owning custom scan must perform. When the plan is constructed, the
% owner removes that marker from the conditions that PostgreSQL would otherwise
% execute. An unowned non-null call reaches a function body that reports an
% unsupported query shape. There is no ordinary scalar alignment fallback.

% \subsection{Inspiration from PostgreSQL extensions}

A full text search interface (\texttt{pg_search}) from ParadeDB was an important motivation for the SQL API design of BioseqDB. 
ParadeDB is an extension written in Rust that uses custom text storage engine under the hood. I deeply inspected its source code to understand how we can implement our own syntax in BioseqDB.
ParadeDB resolves similar problem as FTS can also return multiple hits for a single query.\coderef{paradedb/pg_search/src/postgres/customscan/basescan/projections/score.rs}
{ParadeDB score expression}{score_from_relation}{30}

To return custom results instead of relying on Postgres evaluating the function (We don't want to let Postgres call actual \path{bioseq.score()} function, but rether provide it ourselves. ParadeDB has similar syntax and the same problem)
 ParadeDB replaces score expressions with typed placeholders.
These placefolders are filled from the current search result before projection.
\coderef{paradedb/pg_search/src/postgres/customscan/projections.rs}
  {ParadeDB placeholder replacement}{inject_placeholders}{493}
Similar design was used in BioseqDB - described in Section~\ref{sec:match-placeholder-injection}.

\noindent\begin{minipage}{\linewidth}
\begin{lstlisting}[style=bioseq-search-sql,caption={ParadeDB pg_search examples},label={lst:paradedb-pg-search-examples}]
-- Some examples for pg_search
SELECT id, title, pdb.score(id) AS score
FROM documents
WHERE body ||| 'postgresql search'
ORDER BY score DESC, id;

SELECT id, pdb.snippet(body) AS excerpt
FROM documents
WHERE body &&& 'postgresql indexing';
\end{lstlisting}
\end{minipage}

Listing~\ref{lst:paradedb-pg-search-examples} shows how similar ParadeDB's problem is to my initial API designs.\autocite{pigsty_pg_search}\autocite{neon_pg_search}

\subsection{Paths, plans, and live state}
\label{sec:planner-support-and-ownership}
\label{sec:custom-path-plan-lifecycle}

To implement custom scan, we must understand different objects in Postgres:
\begin{enumerate}
  \item \textbf{Path:} Describes
    possible inputs, estimated work, output row count, and ordering. The
    planner usually compares several paths before selecting one (this is done via genetic query optimizer GQO~\autocite{postgresql_planner_optimizer}\autocite{conway_postgres_optimizer}, but in this case we don't have to know it's internal implementation details, we just provide simple information about path). 
    The path is later converted into a plan. I register serveral own paths to let postgres know how to access the table.
  \item \textbf{Plan:} In our case \texttt{CustomScan} contains child plans, expression trees, and descriptions of the
    search.
  \item \textbf{State:} When plan is executed, postgres creates its state via my code. \texttt{CustomScanState} holds the input cursor, alignment
    buffers for sequences and metadata, and slots for generated hit.
    Postgres will call scan repeatedly to get new rows, so I use the state to hold any information that needs to be preserved between calls.
\end{enumerate}

All types of Postgres API described in table earlier~\ref{tab:pg-scan-extension-points} mainly rely on callbacks (in Postgres called \emph{hooks}). Extension fills the structure with function pointers and let postgres know which call is handled by which function.
In this case PostgreSQL calls \texttt{PlanCustomPath} when it converts a selected custom
path into a plan and later calls \texttt{CreateCustomScanState} to create the
executor state\footnote{PostgreSQL
17, \href{https://www.postgresql.org/docs/17/custom-scan-path.html}
{Creating Custom Scan Paths} and
\href{https://www.postgresql.org/docs/17/custom-scan-plan.html}
{Creating Custom Scan Plans}.}

\begin{figure}[htbp]
\centering
\begin{tikzpicture}[font=\footnotesize]
  \node[anchor=west,font=\bfseries] at (0,4.7) {SQL expressions};
  \node[scan box,minimum width=5.0cm] (project0) at (2.6,3.75)
    {id, score, start};
  \node[scan hit,minimum width=5.0cm] (filter0) at (2.6,2.25)
    {id $>20$ and mmsearch(seq, query)};
  \node[scan input,minimum width=2.7cm] (table0) at (2.6,.75)
    {target rows};
  \draw[scan arrow] (table0) -- (filter0);
  \draw[scan arrow] (filter0) -- (project0);
  \node[anchor=west,font=\bfseries] at (7.0,4.7) {Custom plan};
  \node[scan output,minimum width=4.9cm] (project1) at (9.55,3.75)
    {id and fields of the current hit};
  \node[scan hit,minimum width=4.9cm] (owner) at (9.55,2.25)
    {BioSeqDB custom scan};
  \node[scan input,minimum width=2.7cm] (table1) at (8.05,.75)
    {rows with id $>20$};
  \node[scan input,minimum width=1.6cm] (query1) at (11.45,.75)
    {query};
  \draw[scan arrow] (table1) -- (owner);
  \draw[scan arrow] (query1) -- (owner);
  \draw[scan arrow] (owner) -- (project1);
  \draw[scan arrow,dashed,black!55] (5.45,2.25) -- (6.8,2.25);
  \node[align=center,text=black!65] at (6.1,1.45) {plan\\rewrite};
\end{tikzpicture}
\caption{A conceptual expression tree and the work assigned by the custom
plan. The search becomes the operation that produces hit rows. Ordinary
identifier filtering remains before alignment. The boxes describe logical
roles, rather than a literal \texttt{EXPLAIN} output.}
\label{fig:mmsearch-plan-flow}
\end{figure}

\subsection{What the hooks must establish}

A \emph{hook} is a function that PostgreSQL calls at a defined planning
stage. BioSeqDB installs relation and join path hooks. A relation hook sees
candidate access methods for a table. A join hook sees candidate ways to
combine relations. Their task is more than finding the name
\texttt{mmsearch}:

\begin{itemize}
  \item \textbf{Recognize the operation.} Check the function identity, the
    argument types, the supported expression shape, and the search options.
  \item \textbf{Find its inputs.} Associate each operand with its table alias
    and column, or with a supported expression derived from them. For the
    target operand, this association and its query scope form the
    \emph{target binding} used by result helpers.
  \item \textbf{Propose an owner.} Add a candidate base-scan or supported
    join path whose inputs can supply the search arguments. Selecting a path
    fixes the owner in the resulting plan; the hook does not run alignment.
  \item \textbf{Remove ordinary execution sites.} Remove the owned marker
    from row, vectorized, and storage filter lists. A leftover copy would
    call the scalar failure body after the custom scan had already done the
    search.
\end{itemize}

A planner support function also reports a high cost for scalar
\texttt{mmsearch} evaluation.\footcite{postgresql17_support_functions} Cost
encourages a suitable plan, but does not enforce the search's output rules. The
structural checks and the unowned-call error provide that boundary.
Unsupported join, security, or expression shapes do not become correct merely
because a path has a low estimated cost.
